\begin{proof}[Proof of \cref{obj:lem:inverse-count-arm-risk}]
Take \(C=2^{20}>0\), and fix the parameters, law, and counts in the statement. Throughout, suppress the argument \(P\) in the population quantities. For every cell,
\[
p_j=s_{aj}+s_{1-a,j},\qquad
0\le q_{aj}\le s_{aj},\qquad
0\le\mu_{aj}\le1,\qquad
\sum_{j=1}^{d}p_j=1.
\]
Indeed, the arm probabilities partition the covariate-cell probability, and the outcome-success event is contained in the corresponding arm-cell event. Dividing \(0\le q_{aj}\le s_{aj}\) by \(s_{aj}>0\) gives the conditional-mean bounds; when \(s_{aj}=0\), both \(q_{aj}\) and \(\mu_{aj}\) are zero.

\begin{enumerate}
\item \emph{Exact expectation.}
Define the nonnegative Poisson means
\[
\lambda_j=t s_{aj},\qquad \rho_j=t s_{1-a,j}.
\]
The Poisson first moments give
\[
\mathbb E Z_j=u q_{aj},\qquad \mathbb E W_j=\rho_j.
\]
By \cref{obj:lem:poisson-inverse-moments},
\[
\mathbb E\!\left[\frac1{K_j+1}\right]
=
\begin{cases}
1,&\lambda_j=0,\\[2pt]
\dfrac{1-e^{-\lambda_j}}{\lambda_j},&\lambda_j>0.
\end{cases}
\]
Mutual independence permits factorization of the product expectation:
\[
\begin{aligned}
\mathbb E H_j
&=\frac1u\left(
\mathbb E Z_j+
\mathbb E Z_j\,
\mathbb E\!\left[\frac1{K_j+1}\right]\,
\mathbb E W_j\right)\\
&=q_{aj}\left(1+\rho_j
\begin{cases}
1,&\lambda_j=0,\\[2pt]
\dfrac{1-e^{-\lambda_j}}{\lambda_j},&\lambda_j>0.
\end{cases}\right).
\end{aligned}
\]
The count factors are integrable, the reciprocal factor is bounded by one, and independence makes their products integrable. Thus finite-sum linearity yields
\[
\mathbb E\!\left[\sum_{j=1}^{d}H_j\right]
=
\sum_{j=1}^{d}q_{aj}\left(1+\rho_j
\begin{cases}
1,&\lambda_j=0,\\[2pt]
\dfrac{1-e^{-\lambda_j}}{\lambda_j},&\lambda_j>0.
\end{cases}\right).
\]
% lean: CausalSmith.Stat.AnnotationRarearmFrontier.inverse_count_poisson_arm_mean

\item \emph{Bias bound.}
First suppose \(s_{aj}=0\). Then \(q_{aj}=\mu_{aj}=0\), so
\[
\mathbb E H_j
=p_j\mu_{aj}-s_{1-a,j}\mu_{aj}e^{-t s_{aj}}
=0.
\]
If \(s_{aj}>0\), then \(\lambda_j>0\), and the preceding expectation simplifies to
\[
\begin{aligned}
\mathbb E H_j
&=q_{aj}
+s_{1-a,j}\frac{q_{aj}}{s_{aj}}
  \bigl(1-e^{-t s_{aj}}\bigr)\\
&=(s_{aj}+s_{1-a,j})\mu_{aj}
-s_{1-a,j}\mu_{aj}e^{-t s_{aj}}\\
&=p_j\mu_{aj}-s_{1-a,j}\mu_{aj}e^{-t s_{aj}}.
\end{aligned}
\]
Consequently, in all cells,
\[
\mathbb E\!\left[\sum_{j=1}^{d}H_j\right]-\psi_a
=-\sum_{j=1}^{d}s_{1-a,j}\mu_{aj}e^{-t s_{aj}}.
\]
% lean: CausalSmith.Stat.AnnotationRarearmFrontier.inverse_count_cell_mean_identity

For \(p_j>0\), \cref{obj:ass:overlap} gives
\(s_{1j}\ge\epsilon p_j\) and \(s_{0j}\ge\epsilon p_j\):
multiply its lower propensity bound and its upper propensity bound by
\(p_j\), respectively. For \(p_j=0\), the same inequalities follow from
nonnegativity. Hence, for every \(j\),
\[
s_{aj}\ge\epsilon p_j,\qquad
0\le s_{1-a,j}\mu_{aj}\le s_{1-a,j}\le p_j.
\]
Since every summand in the bias expression is nonnegative,
\[
\left|\mathbb E\!\left[\sum_{j=1}^{d}H_j\right]-\psi_a\right|
=\sum_{j=1}^{d}s_{1-a,j}\mu_{aj}e^{-t s_{aj}}
\le\sum_{j=1}^{d}p_j e^{-t\epsilon p_j}.
\]
The last sum is at most one, because
\[
\sum_{j=1}^{d}p_j e^{-t\epsilon p_j}
\le\sum_{j=1}^{d}p_j=1.
\]
The elementary exponential bound gives
\[
(t\epsilon p_j)e^{-t\epsilon p_j}\le e^{-1},
\qquad
p_j e^{-t\epsilon p_j}\le\frac1{\exp(1)t\epsilon}.
\]
Summing this second bound proves
\[
\left|\mathbb E\!\left[\sum_{j=1}^{d}H_j\right]-\psi_a\right|
\le
\min\left\{1,\frac{d}{\exp(1)t\epsilon}\right\}.
\]
% lean: CausalSmith.Stat.AnnotationRarearmFrontier.inverse_count_mean_bias_bound

\item \emph{Weight moments.}
Use the reciprocal function and numerical constant from
\cref{obj:lem:poisson-inverse-moments}, and define the nonnegative real
random variables
\[
g(K_j)=\frac1{K_j+1},\qquad
C_0=2^{16},\qquad
\omega_j=W_jg(K_j),\qquad
\beta_j=1+\omega_j.
\]
The cited result supplies
\[
\mathbb E[g(K_j)^2]\le\frac{C_0}{(1+\lambda_j)^2},
\qquad
\operatorname{Var}(g(K_j))
\le\frac{C_0\lambda_j}{(1+\lambda_j)^4}.
\]
The usual Poisson moments are
\[
\mathbb E W_j=\rho_j,\qquad
\mathbb E W_j^2=\rho_j^2+\rho_j,\qquad
\operatorname{Var}(W_j)=\rho_j.
\]
Independence of \(W_j\) and \(K_j\) therefore gives
\[
\mathbb E\omega_j^2
=(\rho_j^2+\rho_j)\mathbb E[g(K_j)^2].
\]
In particular, \(\omega_j\) and \(\beta_j\) are square integrable. Since
\((1+\omega_j)^2\le2+2\omega_j^2\), by
\((\omega_j-1)^2\ge0\),
\[
\mathbb E\beta_j^2
\le2+2\mathbb E\omega_j^2
\le2+2C_0\frac{\rho_j^2+\rho_j}{(1+\lambda_j)^2}.
\]
Expanding variance and factoring the independent expectations yields
\[
\begin{aligned}
\operatorname{Var}(\beta_j)
&=\operatorname{Var}(\omega_j)\\
&=(\rho_j^2+\rho_j)\mathbb E[g(K_j)^2]
-\rho_j^2\bigl(\mathbb E g(K_j)\bigr)^2\\
&=\rho_j\mathbb E[g(K_j)^2]
+\rho_j^2\operatorname{Var}(g(K_j))\\
&\le C_0\left\{
\frac{\rho_j}{(1+\lambda_j)^2}
+\frac{\rho_j^2\lambda_j}{(1+\lambda_j)^4}
\right\}.
\end{aligned}
\]
These calculations also cover \(\lambda_j=0\) and \(\rho_j=0\).
% lean: CausalSmith.Stat.AnnotationRarearmFrontier.inverse_count_weight_moments

\item \emph{Cell variance.}
The Poisson moments of the outcome-success count give
\[
\mathbb E(Z_j/u)=q_{aj},\qquad
\mathbb E[(Z_j/u)^2]=q_{aj}^2+\frac{q_{aj}}u,
\qquad
\operatorname{Var}(Z_j/u)=\frac{q_{aj}}u.
\]
Moreover, \(Z_j/u\) is independent of \(\beta_j\), because \(Z_j\) is
independent of the pair \((W_j,K_j)\). Their product is square integrable:
its second moment factors into the two finite second moments. Thus
\[
\begin{aligned}
\operatorname{Var}(H_j)
&=\mathbb E[(Z_j/u)^2]\mathbb E\beta_j^2
-\bigl(\mathbb E(Z_j/u)\bigr)^2
 \bigl(\mathbb E\beta_j\bigr)^2\\
&=\frac{q_{aj}}u\,\mathbb E\beta_j^2
+q_{aj}^2\operatorname{Var}(\beta_j).
\end{aligned}
\]
Define the deterministic cell envelope by
\[
\mathcal V_{\mathrm{cell},j}
=
\frac{q_{aj}}u\left(
2+2C_0\frac{\rho_j^2+\rho_j}{(1+\lambda_j)^2}
\right)
+C_0q_{aj}^2\left(
\frac{\rho_j}{(1+\lambda_j)^2}
+\frac{\rho_j^2\lambda_j}{(1+\lambda_j)^4}
\right).
\]
Multiplying the weight bounds by the nonnegative coefficients
\(q_{aj}/u\) and \(q_{aj}^2\) gives
\[
\operatorname{Var}(H_j)\le\mathcal V_{\mathrm{cell},j}.
\]
% lean: CausalSmith.Stat.AnnotationRarearmFrontier.inverse_count_poisson_cell_variance

\item \emph{Variance of the sum.}
For \(j\ne k\), the count blocks use disjoint members of the mutually
independent collection, so
\[
(Z_j,K_j,W_j)\ \perp\!\!\!\perp\ (Z_k,K_k,W_k),
\qquad
H_j\ \perp\!\!\!\perp\ H_k.
\]
The second assertion follows by applying the measurable function defining
each contribution to its block. Pairwise independence and square
integrability now give
\[
\operatorname{Var}\!\left(\sum_{j=1}^{d}H_j\right)
=\sum_{j=1}^{d}\operatorname{Var}(H_j)
\le\sum_{j=1}^{d}\mathcal V_{\mathrm{cell},j}.
\]
% lean: CausalSmith.Stat.AnnotationRarearmFrontier.inverse_count_cell_blocks_independent

\item \emph{Deterministic cell envelope.}
If \(s_{aj}=0\), then \(q_{aj}=0\), and
\(\epsilon(s_{aj}+s_{1-a,j})\le s_{aj}\) forces
\(s_{1-a,j}=0\). Thus \(p_j=0\) and
\(\mathcal V_{\mathrm{cell},j}=0\).

Suppose henceforth that \(s_{aj}>0\). Overlap implies
\[
p_j\le\frac{s_{aj}}\epsilon\le\frac{p_j}\epsilon,
\qquad
s_{1-a,j}\le\frac{p_j}\epsilon,
\qquad
q_{aj}\le\frac{p_j}\epsilon.
\]
It also gives \(\epsilon s_{1-a,j}\le s_{aj}\), whence
\[
\epsilon s_{1-a,j}^2
\le s_{aj}s_{1-a,j}
\le s_{aj}p_j,
\qquad
\frac{s_{1-a,j}^2}{s_{aj}}\le\frac{p_j}\epsilon.
\]
% lean: CausalSmith.Stat.AnnotationRarearmFrontier.inverse_count_overlap_ratio_bound

The rational factors satisfy
\[
0\le\frac{\lambda_j}{1+\lambda_j}\le1,
\qquad
\lambda_j\le(1+\lambda_j)^2.
\]
Consequently,
\[
\begin{aligned}
\frac{s_{aj}\rho_j^2}{(1+\lambda_j)^2}
&=\frac{s_{1-a,j}^2}{s_{aj}}
  \left(\frac{\lambda_j}{1+\lambda_j}\right)^2
\le\frac{s_{1-a,j}^2}{s_{aj}},\\
\frac{s_{aj}\rho_j}{(1+\lambda_j)^2}
&=s_{1-a,j}\frac{\lambda_j}{(1+\lambda_j)^2}
\le s_{1-a,j},\\
\frac{s_{aj}^2\rho_j}{(1+\lambda_j)^2}
&=\frac{s_{1-a,j}}t
  \left(\frac{\lambda_j}{1+\lambda_j}\right)^2
\le\frac{s_{1-a,j}}t,\\
\frac{s_{aj}^2\rho_j^2\lambda_j}{(1+\lambda_j)^4}
&=\frac{s_{1-a,j}^2}{t s_{aj}}
  \left(\frac{\lambda_j}{1+\lambda_j}\right)^4
\le\frac{s_{1-a,j}^2}{t s_{aj}}.
\end{aligned}
\]
% lean: CausalSmith.Stat.AnnotationRarearmFrontier.inverse_count_rational_factors

Using \(q_{aj}\le s_{aj}\) in the first two bounds yields
\[
\begin{aligned}
q_{aj}\frac{\rho_j^2+\rho_j}{(1+\lambda_j)^2}
&\le
\frac{s_{aj}\rho_j^2}{(1+\lambda_j)^2}
+\frac{s_{aj}\rho_j}{(1+\lambda_j)^2}\\
&\le\frac{s_{1-a,j}^2}{s_{aj}}+s_{1-a,j}
\le\frac{2p_j}\epsilon.
\end{aligned}
\]
Similarly, \(q_{aj}^2\le s_{aj}^2\) and the last two bounds give
\[
\begin{aligned}
q_{aj}^2\left(
\frac{\rho_j}{(1+\lambda_j)^2}
+\frac{\rho_j^2\lambda_j}{(1+\lambda_j)^4}
\right)
&\le
\frac{s_{aj}^2\rho_j}{(1+\lambda_j)^2}
+\frac{s_{aj}^2\rho_j^2\lambda_j}{(1+\lambda_j)^4}\\
&\le\frac{s_{1-a,j}}t
+\frac{s_{1-a,j}^2}{t s_{aj}}
\le\frac{2p_j}{t\epsilon}.
\end{aligned}
\]
Substitution into the cell envelope proves
\[
\begin{aligned}
\mathcal V_{\mathrm{cell},j}
&\le
\frac2u\frac{p_j}\epsilon
+\frac{2C_0}u\frac{2p_j}\epsilon
+C_0\frac{2p_j}{t\epsilon}\\
&=(2+4C_0)\frac{p_j}{u\epsilon}
+2C_0\frac{p_j}{t\epsilon}\\
&\le
(2+4C_0)p_j
\left(\frac1{u\epsilon}+\frac1{t\epsilon}\right).
\end{aligned}
\]
The last inequality adds the nonnegative quantity
\((2+2C_0)p_j/(t\epsilon)\). Together with the zero-mass case, this
establishes the envelope for every cell.
% lean: CausalSmith.Stat.AnnotationRarearmFrontier.inverse_count_cell_variance_envelope

\item \emph{Summation and constant calibration.}
Summing the cell envelopes and using \(\sum_jp_j=1\) gives
\[
\begin{aligned}
\operatorname{Var}\!\left(\sum_{j=1}^{d}H_j\right)
&\le\sum_{j=1}^{d}\mathcal V_{\mathrm{cell},j}\\
&\le(2+4C_0)
\left(\frac1{u\epsilon}+\frac1{t\epsilon}\right).
\end{aligned}
\]
Finally,
\[
2+4C_0=262146\le1048576=2^{20}=C.
\]
Since the reciprocal-intensity sum is nonnegative, this proves
\[
\operatorname{Var}\!\left(\sum_{j=1}^{d}H_j\right)
\le C\left(\frac1{u\epsilon}+\frac1{t\epsilon}\right),
\]
as required.
% lean: CausalSmith.Stat.AnnotationRarearmFrontier.inverse_count_arm_variance_envelope
\end{enumerate}
\end{proof}
